\chapter{Introduction}

\section*{Context and Motivation}

The transformation of video games from a niche pastime into a global cultural and economic phenomenon has fundamentally altered the ways in which individuals interact, compete, and even pursue professional careers. In the era of \textbf{competitive gaming}, players participate in high-stakes environments where success hinges on precision, strategic thinking, and rapid reflexes. Driven by global connectivity and high-speed internet, online multiplayer games have evolved into structured \textbf{competitive ecosystems} featuring ranking systems, seasonal rewards, and financial incentives.

Today, video games function not only as a form of entertainment but also as \textbf{social platforms}, \textbf{economic engines}, and \textbf{professional arenas}. However, this rapid evolution has brought with it a persistent and multifaceted threat: \textbf{cheating}.

Through unauthorized software modifications, certain players seek to gain unfair advantages that disrupt the integrity of competition. These interventions range from memory manipulation and artificial input assistance to hardware-based exploits and network-level interference. As a result, an ongoing \textbf{arms race} has emerged, pitting developers, who design increasingly sophisticated anti-cheat mechanisms, against cheat creators, who constantly innovate to circumvent detection.

Moreover, the deployment of modern anti-cheat technologies raises a series of complex \textbf{legal}, \textbf{ethical}, and \textbf{environmental} concerns. Effective detection often requires deep access to system resources, prompting significant debates around user privacy and digital rights. At the same time, the commercialization of cheat software causes substantial economic harm to developers and challenges regulatory frameworks on a global scale. Additionally, the computational burden imposed by advanced detection mechanisms contributes to a growing environmental footprint within digital infrastructures.

To transcend theoretical discourse, this thesis adopts an applied methodology: the design, development, and evaluation of a custom-built cheat for a widely played multiplayer game. 


\section*{Objectives and Scope}

The primary aim of this thesis is to investigate the phenomenon of cheating in competitive video games through a dual lens: theoretical inquiry and hands-on experimentation. The research is articulated around three core objectives:

\begin{enumerate}
    \item \textbf{To analyze the historical evolution of cheating in video games}, and to assess its influence on game design, player behavior, and the industry's countermeasures.

    \item \textbf{To examine and categorize common cheating techniques}, including software manipulation, visual enhancements, and attacks, in terms of both technical function and competitive impact.

    \item \textbf{To implement and evaluate a custom cheat in a contemporary multiplayer title}, using this case study to test the efficacy of current anti-cheat mechanisms and identify systemic weaknesses.
\end{enumerate}

\section*{Competencies and Alignment with the Master's Program}

This work is directly aligned with the core learning outcomes of the \textbf{Master’s Degree in Cybersecurity}, particularly in the domains of \textbf{reverse engineering}, \textbf{intrusion analysis}, and \textbf{digital forensics}. Throughout the development of the project, the following competencies will be applied and demonstrated:

\noindent\textbf{Basic Competencies}
\begin{itemize}
    \item \textbf{[CB6]} Conduct original research into the security challenges associated with competitive gaming, focusing on cheating and related countermeasures.
    \item \textbf{[CB7]} Apply cybersecurity techniques to the context of video game security, including the detection, development, and circumvention of unauthorized interventions.
    \item \textbf{[CB8]} Evaluate the legal and ethical dimensions of cheat creation and distribution in light of current digital legislation.
\end{itemize}

\noindent\textbf{General Competencies}
\begin{itemize}
    \item \textbf{[CG1]} Analyze real-world intrusion scenarios to identify technical vulnerabilities in online gaming environments.
    \item \textbf{[CG2]} Understand and assess the internal architecture of anti-cheat systems, evaluating their effectiveness and limitations.
    \item \textbf{[CG3]} Produce clear and rigorous technical documentation that reflects expertise in offensive and defensive security techniques.
\end{itemize}

\newpage

\noindent\textbf{Specific Competencies}
\begin{itemize}
    \item \textbf{[CE1]} Understand the psychological and sociological factors that motivate cheating in competitive contexts.
    \item \textbf{[CE2]} Analyze advanced evasion strategies used by cheat developers and their implications for detection models.
    \item \textbf{[CE3]} Monitor and evaluate emerging threats in the field of game security as a subset of broader cybersecurity challenges.
    \item \textbf{[CE4]} Conduct forensic investigations into game memory and process manipulation to detect unauthorized behavior.
    \item \textbf{[CE6]} Critically assess the internal logic and performance trade-offs involved in the implementation of modern anti-cheat technologies.
\end{itemize}

\section*{Document Structure}

The remainder of this thesis is organized as follows:

\begin{itemize}
    \item \textbf{Chapter 1: Introduction} – Presents the motivation, objectives, competencies, and structural overview of the research.
    
    \item \textbf{Chapter 2: From Primitive Survival to Digital Domination} – Explores the human drive for competition and its historical transition into virtual environments, with emphasis on the technological developments that enabled competitive gaming.
    
    \item \textbf{Chapter 3: The Evolution of Cheating: From Debugging Tools to an Organized Industry} – Chronicles the emergence of cheating practices, their evolution into a commercial ecosystem, and the industry's reactive development of anti-cheat systems.
    
    \item \textbf{Chapter 4: The Prevalence and Impact of Cheating in Online Multiplayer Games} – Analyzes the effects of cheating on gameplay integrity, user experience, and industry revenue, with a specific focus on the game \textbf{Counter-Strike 2}.
    
    \item \textbf{Chapter 5: Injection and Detection Techniques} – Provides an overview of the technical mechanisms behind code injection, cheat detection, and the adversarial dynamics between attackers and defenders.
    
    \item \textbf{Chapter 6: Internal Game Structure: Logic Behind Counter-Strike 2} – Investigates the architectural design of CS2's engine, including memory organization, entity management, and reverse engineering techniques.
    
    \item \textbf{Chapter 7: Cheat Features Implemented in Our Cheat for CS2} – Describes the practical implementation of various cheat modules and evaluates their effectiveness against contemporary anti-cheat systems.
\end{itemize}